\documentclass[12pt, a4paper]{article}

\usepackage[utf8]{inputenc}
\usepackage[T1]{fontenc}
\usepackage[russian]{babel}
\usepackage[oglav,spisok,boldsect,eqwhole,figwhole,hyperref,hyperprint,remarks,greekit]{./style/fn2kursstyle}
\graphicspath{{./style/}{./figures/}}

\usepackage{multirow}
\usepackage{supertabular}
\usepackage{multicol}
\usepackage{hyperref}
\usepackage{enumitem}
% Параметры титульного листа
\title{Лабораторная работа 1}
\author{А.\,Р.~Ластра-Грек}
\supervisor{А.\,В.~Череднеченко}
\group{ФН2-51Б}
\date{2023}

% Переопределение команды \vec, чтобы векторы печатались полужирным курсивом
\renewcommand{\vec}[1]{\text{\mathversion{bold}${#1}$}}%{\bi{#1}}
\newcommand\thh[1]{\text{\mathversion{bold}${#1}$}}
%Переопределение команды нумерации перечней: точки заменяются на скобки
\renewcommand{\labelenumi}{\theenumi)}

\usepackage{amsmath}
\usepackage{amssymb}

\usepackage{float, graphicx}

\usepackage{nicefrac}

\usepackage{relsize}
\usepackage{caption}
\usepackage{listings}
\usepackage{xcolor}
\lstset { %
	language=C++,
	%backgroundcolor=\color{black!5}, % set backgroundcolor
	basicstyle=\footnotesize % basic font setting
}

\frenchspacing
\righthyphenmin=2
\linespread{1.25}

\newcounter{controlQCounter}
\setcounter{controlQCounter}{1}

\DeclareMathOperator{\cond}{cond}
\newcommand{\tr}{\ensuremath{\mathrm{T}}}
\newcommand{\bigO}{\ensuremath{\mathrm{O}}}
\newcommand{\linSpan}[2]{\ensuremath{\overline{#1, \, #2}}}
\newcommand{\RR}{\mathbb{R}}
\newcommand{\NN}{\mathbb{N}}

\newcommand{\clock}[4]{\boldsymbol{\mathrm{#1}\!:\!\mathrm{#2} \to \mathrm{#3}\!:\!\mathrm{#4} \blacktriangleright} \quad}
\newcommand{\point}{$\_ \_ \_ \_ \_ \_ \_ \_ \_ \_ \_ \_ \_ \_ \_ \_ \_ \_ \_ \_ \_ \_ \_ \_ \_ $}
\newcommand{\size}{1}

\newenvironment{controlQ}{\textbf{Вопрос \arabic{controlQCounter}.} \addtocounter{controlQCounter}{1} \itshape}{\upshape}
\newenvironment{controlAns}{\textrm{Ответ:} \par }{ \bigskip }

\usepackage{amsthm}

% Создаем стиль оформления по ГОСТ
\newtheoremstyle{gostdef}
{\topsep}     % Отступ сверху
{\topsep}     % Отступ снизу
{\normalfont} % Прямой шрифт текста (по ГОСТ определения не пишут курсивом)
{\parindent}  % Красная строка (абзацный отступ 1.25 см)
{\bfseries}   % Жирный шрифт заголовка
{:}           % Двоеточие после номера
{ 0.5em}      % Пробел после двоеточия перед текстом
{}            % Формат заголовка (по умолчанию)

\theoremstyle{gostdef}
% Команда для нумерации по разделам (Определение 1.1, Определение 1.2) - чаще всего требуют в ВКР
\newtheorem{defenv}{Определение}[section] 

% ИЛИ раскомментируйте строку ниже для сквозной нумерации (Определение 1, Определение 2)
% \newtheorem{defenv}{Определение}

\begin{document}
	
	\begin{center}
		СМОТР
	\end{center}
	
	\begin{enumerate}
	
		\item Слайд 1:
		\item Слайд 2:
		\begin{enumerate}
			\item перейдем к постановке проблемы и целям нашего исследования.
			Как известно, решение двумерных и тем более трехмерных задач газовой динамики сопряжено с огромными вычислительными затратами. Когда мы используем классические схемы высокого разрешения для улавливания ударных волн, такие как схемы Курганова-Тадмора, вычисления на языках высокого уровня, вроде Python, становятся критически медленными.
			\item Поэтому главной целью нашей работы стала разработка собственной, современной и, что самое главное, очень быстрой реализации таких схем для решения систем гиперболических законов сохранения.
			\item Для достижения этой цели мы поставили перед собой три конкретные задачи:
			Во-первых, написать с нуля математическое ядро — двумерный решатель для уравнений Эйлера.
			Во-вторых, радикально ускорить его работу за счет глубокой векторизации и JIT-компиляции на базе фреймворка JAX, чтобы код мог эффективно работать на графических ускорителях.
			И, наконец, в-третьих, строго доказать, что наш быстрый решатель считает правильно. В качестве "золотого стандарта" для проверки точности решения и оценки полученного ускорения мы используем известную открытую библиотеку \texttt{centpy}.
		\end{enumerate}
		\item \textbf{Слайд 3:}
		\begin{enumerate}
			\item Как я уже упоминал, реализация центральных схем высокого разрешения сопряжена с большими вычислительными затратами, особенно при переходе к двумерным задачам. Поэтому для написания программного кода мы отказались от стандартного Python с NumPy и выбрали современный фреймворк JAX от Google.
			\item Первая и, пожалуй, главная причина — это JIT-компиляция, то есть компиляция 'на лету'. В классическом Python циклы по времени работают очень медленно. JAX использует компилятор XLA (Accelerated Linear Algebra), который перед запуском переводит наши функции численного потока в оптимизированный машинный код. Благодаря этому тысячи шагов по времени при решении уравнений выполняются на порядки быстрее.
			\item Второе преимущество — автоматическая векторизация. В двумерном случае нам нужно рассчитывать потоки для каждой ячейки сетки. Вместо того чтобы писать вложенные циклы по координатам X и Y, мы используем встроенную функцию vmap. Она позволяет применить формулу потока сразу ко всей расчетной сетке целиком. Код получается таким же компактным, как для одной ячейки, а математика под капотом обрабатывает весь массив параллельно.
			\item И наконец, нативная поддержка GPU. Зачастую, чтобы перенести расчеты с процессора на видеокарту, нужно переписывать код на C++ и CUDA. Код, написанный на JAX, вообще не требует изменений — он автоматически выполняется на графических ускорителях, если они доступны. Это критически важно для тяжелых задач, вроде двумерного уравнения Эйлера на мелких сетках, которые мы планируем досчитать в ближайшее время.
		\end{enumerate}
		\item \textbf{Слайд 4:}
		\begin{enumerate}
			\item В качестве базовой модели мы используем двумерную систему уравнений Эйлера, которая описывает динамику невязкого сжимаемого газа.
			\item Система записана в консервативной форме, где вектор неизвестных включает плотность, проекции импульса и полную энергию. Система замыкается стандартным уравнением состояния идеального газа с показателем адиабаты $\gamma = 1.4$.
			\item В качестве тестовой задачи мы рассматриваем классическую двумерную задачу Римана. Расчетная область представляет собой единичный квадрат. В начальный момент времени она разбита на четыре равных квадранта, в каждом из которых заданы постоянные значения плотности, вектора скорости и давления.
			\item Граничные условия в задаче — типа Дирихле. На границах мы жестко фиксируем значения, равные начальным параметрам в соответствующих прилегающих квадрантах. Такая конфигурация приводит к возникновению и сложному взаимодействию различных газодинамических разрывов — ударных волн и контактных поверхностей, что является отличным тестом для проверки численной схемы.
		\end{enumerate}
		\item \textbf{Слайд 5:}
		\begin{enumerate}
			\item В классическом методе Годунова на границе каждой ячейки нужно точно решать задачу Римана. Это вычислительно дорого, особенно для сложных уравнений состояния или многомерных задач. Схемы Курганова-Тадмора — это так называемые центральные схемы. Их главное преимущество — они работают как "черный ящик": нам не нужно знать структуру характеристик или вычислять якобианы.
			\item В отличие от более ранних схем (например, Нессяху-Тадмора), схема КТ адаптивно меняет ширину контрольного объема. На гладких участках решения численная вязкость резко падает, а на разрывах — сохраняется ровно в том объеме, чтобы предотвратить осцилляции.
			\item Семейство КТ делится на два лагеря. Полностью дискретные (FD) схемы связывают дискретизацию по времени и пространству в одну сложную формулу. Мы же фокусируемся на полудискретных (SD) схемах. В них мы сначала вычисляем пространственные производные, получая систему обыкновенных дифференциальных уравнений, а затем интегрируем ее по времени любым удобным методом, например, Рунге-Куттой. Это дает большую гибкость.
			\item Разница между 2-м и 3-м порядком (SD2 и SD3) заключается в том, как мы интерполируем данные внутри ячейки: кусочно-линейно для SD2 или кусочно-параболически для SD3. Пока мы реализовали базовую линейную реконструкцию, но модульная архитектура на JAX позволит легко интегрировать SD3 в будущем.
			\item Еще текст для компоновки
			\item Перейдем к математическому аппарату. Для интегрирования системы Эйлера мы выбрали семейство центральных схем Курганова-Тадмора. Их главное преимущество перед upwind-схемами, такими как метод Годунова — это отсутствие необходимости решать сложную и ресурсоемкую задачу Римана на каждой грани ячейки, а также вычислять матрицы Якоби.
			\item При этом, в отличие от классических центральных схем вроде Лакса-Фридрихса, схемы Курганова-Тадмора обладают очень низкой схемной вязкостью. Это достигается за счет точной локальной оценки скорости распространения волн через спектральные радиусы.
			\item В своей работе мы делаем упор на полудискретную формулировку — Semi-Discrete или SD. Она удобна тем, что полностью разделяет пространство и время: по пространству мы используем кусочно-линейную реконструкцию для достижения второго порядка точности, а по времени интегрируем полученную систему ОДУ с помощью строгих TVD-схем Рунге-Кутты
		\end{enumerate}
		\item \textbf{Слайд 6:}
		\begin{enumerate}
			\item Первый шаг в нашей полудискретной схеме — это пространственная реконструкция. Нам нужно перейти от средних значений в центре ячейки к значениям на ее левой и правой границах. Для второго порядка точности мы используем кусочно-линейную реконструкцию, так называемый подход MUSCL.
			\item Однако, по теореме Годунова, любая линейная схема выше первого порядка будет выдавать нефизичные осцилляции, то есть "рябь", на сильных разрывах — например, на ударных волнах. Чтобы этого избежать, схема должна обладать свойством TVD: ее полная вариация не должна возрастать со временем. На практике это достигается введением лимитеров — ограничителей наклона. На слайде приведена формула для производной с использованием лимитера minmod: на гладких участках он сохраняет второй порядок, а вблизи разрывов "обрезает" наклон, локально понижая порядок до первого и подавляя осцилляции
			\item Ниже еще какой то текст нужно все компоновать
			\item На этом слайде показан первый этап нашей схемы — пространственная реконструкция на примере одномерного случая. Наша цель — получить значения функции на левой и правой границах ячейки.
			\item Согласно теореме Годунова, линейные схемы выше первого порядка неизбежно порождают осцилляции на разрывах. Чтобы этого избежать, мы требуем выполнения свойства TVD — общая вариация решения не должна расти со временем, что математически запрещает появление новых локальных экстремумов.
			\item Для этого мы используем ограничители наклона, или лимитеры. На экране приведена формула классического лимитера minmod: на гладких участках он сохраняет второй порядок точности, а на скачках обрезает наклон, предотвращая нефизичную рябь.
			\item Важно отметить, что в нашей программе реализован двумерный случай. Поскольку мы используем декартову сетку, эта же одномерная процедура реконструкции применяется поочередно: сначала вдоль оси X для вычисления потоков по X, а затем вдоль оси Y. Благодаря векторизации в JAX, мы делаем это без дополнительных вычислительных затрат.
			\item Все алгоритмы (реконструкция, вычисление потоков, лимитеры) представлены для одномерного случая ради наглядности, однако в программном коде они полностью реализованы для двумерной постановки с использованием покомпонентного расщепления"			
		\end{enumerate}
		\item \textbf{Слайд 7:}
		\begin{enumerate}
			\item Имея значения $u^+$ и $u^-$ на границах ячеек, мы можем вычислить численные потоки. Формула потока Курганова-Тадмора состоит из двух частей. Первая часть — это простое среднее арифметическое потоков слева и справа от границы. Вторая часть — это диссипативный член, который отвечает за стабильность схемы.
			\item Ключевая особенность схемы KT кроется в коэффициенте $a$. Это максимальная локальная скорость распространения волны на данной границе. Благодаря использованию именно локальной скорости, схема вносит ровно столько искусственной вязкости, сколько нужно для стабильности, не "размазывая" решение.
			\item Подставив потоки в исходное уравнение, мы получаем систему обыкновенных дифференциальных уравнений по времени. Ее мы решаем с помощью метода Рунге-Кутты второго порядка, сохраняющего свойство TVD. Важно отметить, что благодаря библиотеке JAX, расчет потоков по всей сетке векторизован и выполняется одновременно.
		\end{enumerate}
		\item \textbf{Слайд 8:}
		\begin{enumerate}
			\item На экране представлена эволюция поля плотности с течением времени. Решением системы уравнений Эйлера является распределение параметров потока в каждой точке пространства. Цветом закодирована плотность газа. Мы четко видим распад начального разрыва: образование расходящихся ударных волн, областей разрежения и формирование сложных вихревых структур на контактных поверхностях в центре области в результате взаимодействия четырех потоков.
			\item Слева на графике 'А' показаны начальные условия. Расчетная область разделена на четыре зоны, в каждой из которых заданы свои постоянные макроскопические параметры — это плотность, давление и компоненты скорости. На графиках цветовой заливкой показано поле плотности газа.
			\item Справа на графике 'Б' показан результат эволюции системы с течением времени. Из-за распада начальных разрывов на границах квадрантов возникает сложное двумерное взаимодействие волн. Мы четко видим формирование ударных волн, центрированных волн разрежения и контактных поверхностей, которые закручиваются в центре области.
			\item Анализируя этот результат, можно сделать два главных вывода о качестве написанного нами солвера. Во-первых, контуры плотности вблизи ударных волн расположены плотно друг к другу — это значит, что схема обладает низкой численной вязкостью и хорошо 'держит' разрыв. Во-вторых, благодаря использованию TVD-лимитеров, на границах фронтов полностью отсутствуют нефизичные осцилляции, что подтверждает корректность нашей реализации полудискретной схемы Курганова-Тадмора в двумерном случае.
		\end{enumerate}
		\item Слайд 9:
		\begin{enumerate}
			\item Для того чтобы строго доказать, что наш программный комплекс работает корректно и достигает заявленной точности, мы провели анализ сеточной сходимости.
			В качестве тестовой задачи мы взяли эволюцию изэнтропического вихря — это классическое гладкое течение без скачков давления. Так как у нас нет точного аналитического решения, мы применили метод экстраполяции Ричардсона. Суть метода заключается в расчете задачи на трех сгущающихся сетках: крупной, средней и мелкой. Затем мы вычисляем разницу между решениями по интегральной норме $L_1$ и смотрим на скорость убывания ошибки.
			Результаты представлены в таблице слева. Поскольку течение полностью гладкое, нелинейные лимитеры в нашем решателе не активируются и не вносят искусственную диссипацию. Благодаря этому, для всех консервативных переменных мы видим строгий выход на теоретический второй порядок точности. Таким образом, ядро решателя функционирует абсолютно верно
			\item Доказав, что метод работает на втором порядке, мы перешли к более сложной и реальной физике — двумерной задаче Римана, которую мы показывали на предыдущих слайдах. Главная особенность этого течения — наличие сильных газодинамических разрывов, таких как ударные волны.
			Мы провели аналогичный анализ Ричардсона на сетках от 50 до 200 узлов. В таблице видно, что глобальный порядок точности упал с двойки до значений порядка 0.6–0.8.
			Это не ошибка программной реализации, а фундаментальное математическое свойство, описанное в теореме Годунова. Невозможно построить линейную схему высокого порядка, которая бы не давала нефизичных осцилляций на ударных волнах. Чтобы предотвратить эти "осцилляции", в нашем коде срабатывают нелинейные лимитеры. Они распознают разрыв и намеренно вносят вязкость, локально снижая точность схемы до первого порядка. Поскольку норма $L_1$ "собирает" ошибку по всей области, наличие таких зон с первым порядком закономерно снижает общую интегральную оценку сходимости
		\end{enumerate}
		\item Слайд 10:
		\item Слайд 11:
		\begin{enumerate}
			\item На этом слайде представлен один из главных результатов нашей работы — сравнительный анализ производительности.
			\item Мы замерили время, необходимое для решения двумерной задачи Римана до момента времени $t = 0.4$, в зависимости от размера пространственной сетки. Синяя кривая — это эталонная открытая библиотека \texttt{centpy}, написанная на классическом Python. Красная кривая — наша собственная реализация на базе фреймворка JAX.
			\item Как видно из графика, на очень грубых сетках, например 40 на 40 ячеек, время работы решателей примерно одинаково. Это связано с тем, что JAX тратит первоначальное время на JIT-компиляцию вычислительного графа.
			\item Однако по мере измельчения сетки ситуация кардинально меняется. Алгоритм на чистом Python начинает "задыхаться" в циклах, и время его работы стремительно улетает вверх, достигая более двух минут на сетке 320 на 320 ячеек. В то же время наша реализация на JAX благодаря глубокой векторизации тензорных операций масштабируется практически линейно. На той же самой мелкой сетке 320 на 320 наш решатель выполняет работу всего за четыре с половиной секунды.
			\item Таким образом, за счет применения современных подходов к аппаратному ускорению вычислений, мы добились колоссального ускорения программы — более чем в 28 раз на плотных сетках, при сохранении полной математической строгости метода
		\end{enumerate}
	\end{enumerate}
		

\end{document}